\documentclass[titlepage,a4paper]{jsarticle}
\usepackage{../../../sty/import}% 各種パッケージインポート
\usepackage{../../../sty/title}% タイトルページの変更

%% タイトルページの変数
% レポートタイトル
\title{原子力材料と核燃料レポート}
% 提出日
\expdate{\today}
% 科目名
\subject{原子力材料と核燃料}
% 分野
\class{情報経営システム工学分野}
% 学年
\grade{M1}
% 学籍番号
\mynumber{24336488}
% 記述者
\author{本間三暉}

\begin{document}
% titleページ作成
\maketitle
\section{核反応A(a,b)Bを説明せよ．}
核反応とは，入射粒子が標的となる原子核と衝突し，原子核に変化が生じる反応である．
核反応は一般に，
\[
A + a \rightarrow B + b
\]
または，
\[
A(a,b)B
\]
と表される．

ここで，$A$ は標的核，$a$ は入射粒子，$B$ は生成核，$b$ は生成粒子を表す．
入射粒子や生成粒子には，中性子 $n$，陽子 $p$，$\alpha$ 粒子，$\gamma$ 線などがある．
また，生成粒子 $b$ は一つとは限らず，複数個生じる場合もある．

例えば，
\[
{}^{14}_{7}\mathrm{N}(\alpha,p){}^{17}_{8}\mathrm{O}
\]
という反応は，
\[
{}^{14}_{7}\mathrm{N} + {}^{4}_{2}\mathrm{He}
\rightarrow
{}^{1}_{1}\mathrm{H} + {}^{17}_{8}\mathrm{O}
\]
を意味する．
この反応では，窒素14を標的核 $A$ とし，$\alpha$ 粒子を入射粒子 $a$ として衝突させた結果，陽子 $p$ が生成粒子 $b$ として放出され，酸素17が生成核 $B$ として生じる．
反応前後で質量数と原子番号を比較すると，質量数は $14+4=1+17=18$，原子番号は $7+2=1+8=9$ となり，保存されていることが分かる．

このように，$A(a,b)B$ という表記は，標的核にどのような粒子を入射し，どのような粒子が放出され，どのような生成核ができるかを簡潔に示す記法である．
さらに，生成核 $B$ が不安定な場合には，その後 $\alpha$ 線，$\beta$ 線，$\gamma$ 線などを放出してより安定な核種へ変化することがあり，このような放射性核種が生成されることを放射化という．

\begin{thebibliography}{99}
\bibitem{講義資料}原子力材料と核燃料 講義資料
\end{thebibliography}

\end{document}